\documentclass[11pt]{amsart}
\usepackage{amsmath,amssymb,amsthm}
\usepackage[margin=1.2in]{geometry}
\usepackage{hyperref}

\newtheorem{theorem}{Theorem}[section]
\newtheorem{lemma}[theorem]{Lemma}
\newtheorem{proposition}[theorem]{Proposition}
\newtheorem{corollary}[theorem]{Corollary}
\theoremstyle{definition}
\newtheorem{remark}[theorem]{Remark}

\newcommand{\Z}{\mathbb{Z}}
\newcommand{\Q}{\mathbb{Q}}
\newcommand{\im}{\operatorname{Im}}
\newcommand{\re}{\operatorname{Re}}

\title{The additive ladder: congrua sets, Gaussian integers,
and partial impossibility theorems for $3\times3$ magic squares of squares}
\author{CosmicWill}
\date{Draft of August 30, 2026 --- comments welcome}

\begin{document}

\begin{abstract}
Whether a $3\times3$ magic square with nine distinct square entries
exists is a well-known open problem. Any such square has center
$m^2$ and row/column offsets lying in the \emph{congrua set}
$D(m) = \{2ef : e^2+f^2=m^2\}$, and its corners are square if and
only if $D(m)$ contains an additive quadruple $U, V, U{+}V, U{-}V$.
We prove that $D(m)$ admits \emph{no additive relation whatsoever}
--- no solution of $\varepsilon_1 d_1 + \varepsilon_2 d_2 +
\varepsilon_3 d_3 = 0$ --- whenever the split part of $m$ (its
largest divisor supported on primes $\equiv 1 \bmod 4$) is $p^a$,
$pq$, or $p^2q$. Consequently \emph{the center of any $3\times3$
magic square of squares is divisible by three distinct primes
$\equiv 1 \bmod 4$, or its split part is of the form $p^3q$ or
$p^2q^2$.} The proofs pass through a reformulation of additive
relations as six-term vanishing sums on a norm torus in $\Z[i]$,
and are ultimately carried by elementary machinery: valuations in
free multiplicative groups, an exact tangent-half-angle
factorization calculus, divisibility trees, and the classical
quartic descents of Fermat --- both $x^4-y^4=z^2$ and $x^4+y^4=z^2$
occur as terminal cases, alongside the non-congruence of $2$ and
$3$. Every enumeration, identity, and case in the proof is
machine-verified; the verification suite and all artifacts are
public.
\end{abstract}

\maketitle

\section{Introduction}

A \emph{magic square of squares} (MSS3) is a $3\times3$ magic
square whose nine entries are distinct perfect squares. Whether one
exists is open. Classical analysis reduces the problem to
arithmetic of \emph{congrua}: the entries are
$m^2 \pm U$, $m^2 \pm V$, $m^2 \pm U \pm V$ and $m^2$, and the
row/column conditions force $U, V \in D(m)$ where
\[
D(m) \;=\; \{\, 2ef \;:\; e^2 + f^2 = m^2,\ e, f > 0 \,\}.
\]
The corner entries are squares if and only if additionally
$U + V \in D(m)$ and $U - V \in D(m)$: an \emph{additive
quadruple} inside a single congrua set. In particular an MSS3
yields an \emph{additive triple} $d_1 + d_2 = d_3$ in $D(m)$, so
the following conjecture implies that no MSS3 exists.

\begin{center}
\emph{Conjecture (A3.C). No congrua set $D(m)$ contains an
additive triple.}
\end{center}

Exhaustive computation verifies A3.C for all $m \le 10^7$. This
paper proves it unconditionally on three infinite families,
organized by the \emph{split part} of $m$: writing
$m = 2^s r \prod p_i^{a_i}$ with $r$ supported on primes
$\equiv 3 \bmod 4$ and $p_i \equiv 1 \bmod 4$ distinct, the split
part is $\prod p_i^{a_i}$. Only the split part matters: $D(m)$ is
a scalar multiple of the congrua set of the split part.

\begin{theorem}[$\omega = 1$]\label{thm:omega1}
If the split part of $m$ is a prime power $p^a$, then no relation
$\varepsilon_1 d_1 + \varepsilon_2 d_2 + \varepsilon_3 d_3 = 0$
holds with $d_i \in D(m)$, $\varepsilon_i \in \{\pm1\}$
(repetitions allowed).
\end{theorem}

\begin{theorem}[$pq$]\label{thm:pq}
The same holds if the split part is $pq$ with $p \ne q$.
\end{theorem}

\begin{theorem}[$p^2q$]\label{thm:p2q}
The same holds if the split part is $p^2 q$.
\end{theorem}

\begin{corollary}\label{cor:mss3}
The center of any $3\times3$ magic square of squares is divisible
by three distinct primes $\equiv 1 \bmod 4$, or has split part of
the form $p^3q$ or $p^2q^2$ with $p, q \equiv 1 \bmod 4$.
\end{corollary}

The proofs share a five-layer architecture which appears to close
every exponent box $(a,b)$ it has been applied to, with residual
case families visibly parametric in $(a,b)$; this suggests a
uniform theorem for all $m$ with at most two distinct split primes
(\S\ref{sec:outlook}).

\subsection*{Method} Additive relations in $D(m)$ are exactly
six-term vanishing sums
$\sum_j \varepsilon_j (w_j - m^4/w_j) = 0$
with $w_j = z_j^2$ on the norm-$m^4$ torus of $\Z[i]$
(\S\ref{sec:setup}). For split part $p^aq^b$ the $w_j$ are
monomials $\sigma^j\tau^k$ in a free rank-two subgroup of
$\Q(i)^\times$, and a relation is a vanishing sum of at most six
monomials. The layers:
(1) \emph{valuation pruning} in the free group kills most sign
patterns outright;
(2) an exact \emph{tangent-half-angle calculus} --- $\sigma =
(1+it_1)/(1-it_1)$ with $t_1$ \emph{rational} --- converts each
surviving pattern into an integer polynomial identity which, for
coherent patterns, factors completely into terms that would
force $\sigma^\alpha\tau^\beta = \pm1$, impossible in a free
group;
(3) finite \emph{congruence sieves} under Pythagorean side
conditions;
(4) \emph{Gaussian collapse}: identities such as
$p^{2d} \pm \ell^{2d} = \ell^d(\bar\ell^d \pm \ell^d)$
($\ell = \lambda^2$, $\lambda$ the Gaussian prime over $p$)
compress the remaining relations to statements a $q$-adic
valuation kills;
(5) \emph{divisibility trees} whose terminal cases are classical:
the two Fermat quartic theorems, the non-congruence of $2$ and
$3$, and short explicit descents.

\subsection*{Verification} Every enumeration and every identity in
this paper is machine-verified in an open suite (148 checks at the
time of writing); \S\ref{sec:methodology} describes the
methodology, including an enumeration-completeness audit performed
after a gap --- omission of the all-equal sign class --- was caught
by the suite's own consistency checks and repaired. We regard
enumeration completeness as a theorem obligation, and prove it by
comparison against an enumeration-independent brute construction.

\section{Setup and the $\Z[i]$ reformulation}\label{sec:setup}

With $z = e + if$, $\im(z^2) = 2ef$, so
$D(m) = \{\,|\im(z^2)| : z \in \Z[i],\ |z|^2 = m^2\,\}$.
Writing $w_j = z_j^2$, so that $|w_j| = m^2$ and $\bar w_j =
m^4/w_j$, a signed additive relation is exactly the six-term
vanishing sum $\sum_j \varepsilon_j (w_j - m^4/w_j) = 0$.
For split part $p^aq^b$, unique factorization in $\Z[i]$ gives
\[
D(m) = \bigl\{\, m^2\,\bigl|\im(\sigma^j\tau^k)\bigr| \;:\;
(j,k) \in [-a,a]\times[-b,b] \setminus \{0\} \,\bigr\},
\qquad
\sigma = \frac{\lambda^4}{p^2},\quad \tau = \frac{\mu^4}{q^2},
\]
where $\lambda, \mu$ are Gaussian primes over $p, q$. The
valuations $v_{\lambda,\bar\lambda,\mu,\bar\mu}(\sigma^j\tau^k) =
(2j, -2j, 2k, -2k)$ show $\langle\sigma,\tau\rangle$ is free of
rank two and no monomial is $\pm1$.

\begin{lemma}[equal-modulus rigidity]\label{lem:rigidity}
No three nonzero elements of $\Q(i)$ of equal absolute value
satisfy $\pm a \pm b \pm c = 0$.
\end{lemma}
\begin{proof}
Absorbing signs, $a+b=c$; dividing by $c$ gives $\alpha + \beta =
1$ with $|\alpha| = |\beta| = 1$ in $\Q(i)$, whence
$\re\,\alpha = \tfrac12$ and $\alpha = \zeta_6^{\pm1} \notin
\Q(i)$.
\end{proof}

\begin{proposition}[degenerate subsums]\label{prop:degenerate}
Every vanishing proper subsum of the six-term sum has size $2$ or
$4$, and a size-$2$ vanishing forces equal congrua. Hence genuine
additive triples of distinct congrua give nondegenerate vanishing
sums.
\end{proposition}
\begin{proof}
Sizes $1, 5$ are impossible (all terms have modulus $m^2$), size
$3$ by Lemma~\ref{lem:rigidity}; a size-$2$ relation forces $w'
\in \{\pm w, \pm\bar w\}$, giving $|\im w'| = |\im w|$.
\end{proof}

\section{One split prime: proof of Theorem \ref{thm:omega1}}

With a single split prime, $D(m) = \{\, m^2 |\im \sigma^k| : 1 \le
k \le a \,\}$, where $\sigma = \lambda^4/p^2$ satisfies the
primitive integer polynomial $R(x) = p^2x^2 - 2Cx + p^2$ with
$C = \re \lambda^4$ and $\gcd(C, p) = 1$ (if $\lambda \mid 2C =
\lambda^4 + \bar\lambda^4$ then $\lambda \mid \bar\lambda^4$,
contradicting coprimality). A signed relation becomes, after
absorbing orientations,
$\sum_i \varepsilon_i'(\sigma^{k_i} - \sigma^{-k_i}) = 0$.
If all $k_i$ are equal this reads (odd integer) $\cdot\,
2i\im\sigma^K = 0$ with $\im \sigma^K \neq 0$: impossible.
Otherwise multiplying by $\sigma^{K}$, $K = \max k_i$, exhibits
$\sigma$ as a root of
\[
Q(x) = \sum_i \varepsilon_i'\bigl(x^{K+k_i} - x^{K-k_i}\bigr),
\]
an integer polynomial with coefficients in $\{0,\pm1,\pm2,\pm3\}$
and $\mathrm{lc}(Q) = -Q(0) \ne 0$. Since $\sigma$ is nonreal, its
minimal polynomial over $\Q$ is $x^2 - (2C/p^2)x + 1$, with
primitive integer form $R$. By Gauss's lemma $Q = \gamma R T$ with
$T \in \Z[x]$ and $\gamma = \operatorname{cont}(Q) \le 3$;
comparing leading coefficients,
$|\mathrm{lc}(Q)| \le 3 < p^2 \le |\gamma\, p^2\, \mathrm{lc}(T)|$
--- impossible for $p \ge 5$. \qed

\section{Two split primes I: the framework and the $(1,1)$ box}
\label{sec:pq}

% TODO(next session): full text of the valuation prune lemma, the
% tan-half factorization lemma with the zero dictionary, the
% complete Family I/II/III trees, and Lemmas L1--L5 with the
% self-contained descents (import from the repository document
% docs/attacks/A3-simultaneous-congrua.md, section 2.6).

\section{Two split primes II: the $(2,1)$ box}\label{sec:p2q}

% TODO(next session): census methodology; the Gaussian collapse
% lemma; the beta-family trees; the sliver descent (leg
% decomposition against the q < p^2 window); the E3-minus descent
% ending in Fermat's x^4 + y^4 = z^2; the closure ledger.

\section{Machine verification and the completeness audit}
\label{sec:methodology}

% TODO(next session): suite description; the enumeration
% completeness theorem (brute construction equals corrected
% enumeration, both boxes, exact set equality: 16 + 24 and
% 140 + 84 canonical patterns); the gap-and-repair episode as a
% case study; artifact inventory and reproduction instructions.

\section{Outlook}\label{sec:outlook}

The residual case families of every box analyzed so far are
parametric in the exponent box, and the five-layer machinery is
uniform; we conjecture and are pursuing a single theorem for all
$m$ with at most two distinct split primes. The known-$m$ frontier
of the MSS3 problem itself now lies at three or more distinct
split primes, where the monomial group has rank three; whether the
present methods extend is the natural next question.
% TODO: the convergent-sum heuristic; the omega >= 3 census.

\end{document}
